\documentclass[11pt]{article}
\usepackage[T1]{fontenc}
\usepackage{lmodern,amsmath,amssymb,amsthm,booktabs}
\usepackage[margin=1in]{geometry}
\usepackage[colorlinks=true,urlcolor=blue,citecolor=blue]{hyperref}
\newtheorem{theorem}{Theorem}
\newtheorem{lemma}[theorem]{Lemma}
\newtheorem{corollary}[theorem]{Corollary}
\newtheorem{conjecture}[theorem]{Conjecture}
\theoremstyle{remark}
\newtheorem{remark}[theorem]{Remark}
\newcommand{\rem}{\operatorname{rem}}
\newcommand{\eq}{\operatorname{eq}}
\emergencystretch=2em
\title{Fixed programs have few shapes:\\a reduction of the Main Conjecture to a counting problem}
\author{Research note accompanying the uniform-process construction}
\date{October 20, 2026}
\begin{document}
\maketitle

\begin{abstract}
A computation from the remainder function determines a \emph{shape}: the
straight-line sequence of remainder instructions it performs, each operand
named by where it came from. We show that a fixed recursive program,
deterministic or nondeterministic, produces at most $C_P^{\,c+1}$ shapes
among its computations with at most $c$ remainder calls. Unrestricted derivations have about
$c^{2c}$ shapes. Under the hypothesis below, this difference accounts for the
whole gap between $\log a$ and $\log a/\log\log a$. Both open questions left by the companion manuscript
\cite{companion} would be settled by a single counting hypothesis: how many pairs
$(a,b)\leq N$ can one shape of length $d$ bring to the value $1$? If no shape
of length at most $c_0\log N$ succeeds on more than $N^{2-\delta}$ pairs, then
every fixed program makes at least $r_P\log a$ remainder calls on almost all
coprime pairs. That is the Main Conjecture, in a strong form. The same
hypothesis also makes the companion bound $O(\log a/\log\log a)$ tight for
almost all coprime pairs. Without any hypothesis, a single program with $o(\log a)$
calls would force shapes of length $o(\log N)$ that succeed on $N^{2-o(1)}$
pairs. The Shape-Count Lemma and both reductions are formalized in Lean; the
counting hypothesis is open.
\end{abstract}

\section{Shapes}

In this note $\log$ is the natural logarithm, as in the Lean statements. The
companion manuscript uses base $2$; this rescales only the constants $c_0$ and
$r_P$ below, and $\log N/\log\log N$ is asymptotically independent of the base.

A \emph{shape} of length $d$ is a sequence $S=((i_1,j_1),\ldots,(i_d,j_d))$
with $i_t,j_t\in\{0,1,\ldots,t\}$. On an input $(a,b)$ put $v_0=a$, $v_1=b$
and $v_{t+1}=\rem(v_{i_t},v_{j_t})$. The shape is \emph{valid} on $(a,b)$ if
every divisor $v_{j_t}$ is nonzero. It \emph{reaches $1$} if moreover some
$v_{t+1}=1$. Thus a remainder-only derivation in the sense of
\cite{companion} is exactly a valid shape. Let
\[
P_S(N)=\#\{(a,b):2\leq b<a\leq N,\ S\text{ is valid on }(a,b)\text{ and reaches }1\}.
\]
There are $\prod_{t=1}^{d}(t+1)^2=((d+1)!)^2$ shapes of length $d$. This is
the count that underlies $\log a/\log\log a$.

\section{The shape count of a fixed program}

\subsection{Programs and computations}\label{sec:model}
We make the operational model explicit. It is the standard call-by-value
reading of the recursive programs of \cite[\S2A, \S2E]{book}, and we take it
as the definition of a computation and of its cost.

\paragraph{Structure.} The structure is $(\mathbb N,\rem,\eq_0)$ or
$(\mathbb N,\rem,\eq_0,\eq_1)$. There are two sorts, $\iota$ (natural numbers)
and $\mathbf b$ (truth values $\mathrm{tt},\mathrm{ff}$). $\rem(x,y)$ is
defined only for $y>0$, and $\eq_0(x)$, $\eq_1(x)$ test $x=0$ and $x=1$. The
structure has no constants of sort $\iota$, so a program has none either.
This matters: a constant $1$ would make every shape question trivial.

\paragraph{Programs.} A program $P$ has finitely many function symbols
$p$, each with an arity $k_p$, argument sorts and a result sort. It has
finitely many equations $p(x_1,\ldots,x_{k_p})=e$, possibly several for the
same $p$ (nondeterminism), and a main term $e_0$ in the input variables
$x_1,x_2$. Terms are built from variables, the constants
$\mathrm{tt},\mathrm{ff}$, the primitives $\rem(s,t)$, $\eq_0(s)$, $\eq_1(s)$,
conditionals $\mathrm{if}\ s\ \mathrm{then}\ t\ \mathrm{else}\ u$ with $s$ of
sort $\mathbf b$, and calls $p(s_1,\ldots,s_{k_p})$, all well sorted. The
variables of a body are among its equation's parameters. Let $\mathcal E$ be
the finite set of subterms of bodies and of $e_0$. Let $K\geq2$ bound the
arities, and let $A$ bound the number of immediate subterms of a term in
$\mathcal E$. A call has its arguments as immediate subterms, plus one more
child, the body.

\paragraph{Computations.} An \emph{environment} assigns values to the
variables of an equation (or of $e_0$). A \emph{computation tree} is a finite
tree whose nodes are triples (environment $\rho$, term $e\in\mathcal E$, value
$v$), built by these rules.
\begin{itemize}
\item \textbf{Variables and constants:} a leaf with $v=\rho(x)$, or $v$ the constant.
\item \textbf{Primitives:} for $\rem(s,t)$, $\eq_0(s)$, $\eq_1(s)$ the children
are computations of $s$ (and $t$) in $\rho$ with values $x$ (and $y$). For
$\rem$ we need $y>0$, and then $v=\rem(x,y)$. For the tests, $v$ is the truth
value of $x=0$ or $x=1$.
\item \textbf{Conditionals:} the children are a computation of $s$ in $\rho$
with value $\beta$, and a computation of $t$ (if $\beta=\mathrm{tt}$) or of $u$
(if $\beta=\mathrm{ff}$) in $\rho$ whose value is $v$.
\item \textbf{Calls:} for $p(s_1,\ldots,s_k)$ the children are computations of
each $s_i$ in $\rho$ with values $w_i$, then a computation of the body $e$ of
\emph{some} equation for $p$, in the environment $x_i\mapsto w_i$, whose value
is $v$.
\end{itemize}
A \emph{convergent computation} of $P$ on $(a,b)$ is a computation tree for
$e_0$ in $x_1\mapsto a$, $x_2\mapsto b$. Its value is the value at the root,
and its \emph{cost} is its number of $\rem$ nodes.
\begin{itemize}
\item \textbf{Correctness:} following the book's convention for nondeterministic
  programs, $P$ \emph{computes} $f$ if every convergent computation has the
  correct value and one exists on every input.
\item \textbf{Complexity:} $c_P(a,b)$ is the least cost of a convergent computation on
  $(a,b)$.
\end{itemize}

Divergence never enters: a convergent computation is a finite tree, so any
amount of free, cost-zero recursion in it is finite. Executing the tree in
post-order is what the recursive machine does.

\subsection{Provenance and shapes}
Fix a convergent computation $\mathcal T$ on $(a,b)$. List its $\rem$ nodes
in post-order as $n_1,\ldots,n_c$. Every $\iota$-valued node $z$ gets a
\emph{provenance} $\pi(z)\in\{0,\ldots,c+1\}$, and every environment entry of
sort $\iota$ gets one too, by induction along the tree:
\begin{itemize}
\item the root environment gives $x_1$ provenance $0$ and $x_2$ provenance $1$;
\item a variable leaf has the provenance of its environment entry;
\item the node $n_t$ has provenance $t+1$;
\item a conditional has the provenance of its chosen branch;
\item a call has the provenance of its body child;
\item the body's environment entry $x_i$ has the provenance of the $i$-th
argument child.
\end{itemize}
The \emph{shape} of $\mathcal T$ is
$S(\mathcal T)=((\pi(\ell_1),\pi(r_1)),\ldots,(\pi(\ell_c),\pi(r_c)))$, where
$\ell_t,r_t$ are the two children of $n_t$. By induction along the tree,
every provenance carried by a node or an environment entry is $0$, $1$, or
$s+1$ for a $\rem$ node $n_s$ that is completed before that node or entry is
used: a body's environment entries come from argument children, which precede
the body in post-order. The operands $\ell_t,r_t$ are completed before $n_t$,
so $\pi(\ell_t),\pi(r_t)\leq t$, and this is a shape.

\begin{lemma}[Soundness]\label{lem:sound}
Run $S(\mathcal T)$ on $(a,b)$ with values $v_0,\ldots,v_{c+1}$. Every
$\iota$-valued node $z$ has value $v_{\pi(z)}$, and the same holds for every
$\iota$-valued environment entry. In particular $S(\mathcal T)$ is valid on
$(a,b)$, and every $\iota$-value occurring in $\mathcal T$ is one of the
$v_i$.
\end{lemma}

\begin{proof}
Induction in post-order. Each rule for $\pi$ copies a value together with its
provenance. At $n_t$ the children have values $v_{\pi(\ell_t)}$ and
$v_{\pi(r_t)}$, and the latter is positive because the $\rem$ rule requires
it. So the node's value is $\rem(v_{\pi(\ell_t)},v_{\pi(r_t)})=v_{t+1}$.
\end{proof}

\begin{lemma}[Light subtrees select]\label{lem:light}
Call a node \emph{light} if its subtree contains no $\rem$ node. A light
$\iota$-valued node $z$ evaluated in environment $\rho$ has
$\pi(z)=\pi(\rho(x))$ for some variable $x$ of $\rho$.
\end{lemma}

\begin{proof}
Induction on the subtree. A variable leaf is immediate. A light node cannot
be a $\rem$ node, and no constant has sort $\iota$. A conditional inherits
from its chosen branch, which is evaluated in $\rho$. A call inherits from its
body, which is evaluated in an environment whose entries are the arguments.
The arguments are light $\iota$-valued nodes in $\rho$, so by induction each
has the provenance of a variable of $\rho$.
\end{proof}

\subsection{The Shape-Count Lemma}

\begin{lemma}[Shape-Count Lemma]\label{lem:count}
There is a constant $\Sigma_P$, depending only on $P$, such that the shapes of
convergent computations of $P$ with cost $c$, on all inputs together, number
at most $(2\Sigma_P)^{4c+2}$. Hence those of cost at most $c$ number at most
$C_P^{\,c+1}$, with $C_P=(2\Sigma_P)^5$.
\end{lemma}

\begin{proof}
Let $\mathcal T$ have cost $c$. Call a node \emph{heavy} if it is not light.
Every ancestor of a heavy node is heavy, so for $c\geq1$ the heavy nodes form
a subtree $H$ containing the root. Its leaves are $\rem$ nodes. Fewer than
$c$ of its nodes have two or more heavy children (``branching'' nodes),
because a tree has fewer branching nodes than leaves. Call every other
non-$\rem$ node of $H$ a \emph{chain node}: it has exactly one heavy child.

\emph{The data at a node.} For a node $z$ of $H$ evaluated in environment
$\rho_z$ with variables $x_1,\ldots,x_k$ ($k\leq K$), let the \emph{label} of
$z$ record:
\begin{itemize}
\item the term $e_z\in\mathcal E$;
\item for a call, which equation was used;
\item for each child of $z$ (at most $A+1$ of them), whether it is heavy, and
if it is light and $\iota$-valued, an index $i\leq k$ with
$\pi(\text{child})=\pi(\rho_z(x_i))$ (Lemma~\ref{lem:light}). For the body
child of a call, which lives in the new environment, record instead which
argument position supplies its provenance.
\end{itemize}
Given the provenances of $\rho_z$ and of the results of the heavy children of
$z$, the label determines the provenances of the following, by the rules
defining $\pi$:
\begin{itemize}
\item the environments of all children of $z$;
\item the result of $z$;
\item for a $\rem$ node, its operand pair.
\end{itemize}
There are at most $\Lambda_P=|\mathcal E|\cdot(\#\text{equations})\cdot
(2(K+1))^{A+1}$ labels.

\emph{Contracting chains.} Take a maximal path $z_0\to z_1\to\cdots\to z_m$ in
$H$ whose interior nodes $z_1,\ldots,z_{m-1}$ are chain nodes. All siblings
along it are light, so by induction along the path two things hold:
\begin{itemize}
\item every $\iota$-entry of $\rho_{z_m}$ has the provenance of an entry of
$\rho_{z_1}$, which gives a map $\phi\colon\{1..K\}\to\{1..K\}$;
\item the result of $z_1$ has the provenance either of the result of $z_m$ or
of an entry $\rho_{z_1}(x_i)$, or it is a truth value. That is a datum
$\psi\in\{1..K\}\cup\{\mathrm{bottom},\mathrm{bool}\}$.
\end{itemize}
Contract the path to an edge labelled $(\phi,\psi)$. There are at most
$K^K(K+2)$ edge labels. The contracted tree $H'$ has the following nodes: the
$c$ $\rem$ nodes, fewer than $c$ branching nodes, and the root of $H$ if it
is a chain node. So it has at most $2c$ nodes and at most $2c$ edges.

\emph{Encoding.} Write $H'$ in preorder, each node written as its label and
followed by the labels of the edges to its children, with the children in
their order in $\mathcal T$. This gives a word of length at most
$2c+2c\leq4c+2$ over an alphabet of size $\Sigma_P=\Lambda_P+K^K(K+2)$.

The word determines $S(\mathcal T)$. Decode recursively in evaluation order,
starting at the root with provenances $0,1$ for $x_1,x_2$: given the
provenances of $\rho_z$, treat the children of $z$ in order. The label of $z$
gives the environment provenances of each child from those of $\rho_z$ and,
for a body, from the results of the argument children already decoded. An
edge label transports environment provenances down a contracted chain and the
result provenance back up. Decoding a heavy child returns its result
provenance, and the label then gives that of $z$. So every $\rem$ node's
operand pair is determined. The post-order of the $\rem$ nodes is also
determined, since light subtrees contain none. For $c=0$ the shape is empty.
Finally, there are at most $\sum_{\ell\leq4c+2}\Sigma_P^\ell\leq(2\Sigma_P)^{4c+2}$
words of length at most $4c+2$.
\end{proof}

The proof uses only finiteness of $\mathcal E$, bounded arity, and the
absence of $\iota$-constants; the Lean version below shows that bounded arity
alone suffices, with a weaker constant. Unbounded cost-zero recursion, truth values,
conditionals and nondeterministic choice are all absorbed into the labels.

\begin{remark}[Register programs]\label{rem:width}
The constant is crude, but its order can be pinned down for register
programs. Consider runs of width $W$ whose steps overwrite a register by a
remainder of two registers, route registers (any map of registers to
registers), or run a subcomputation on a copy of the registers and overwrite
one register by its result; the program of \cite{fixedwidth} is of this
kind. We sketch why such runs have at most $(8W^3)^{5c+1}$ shapes with $c$
remainder calls; this sketch is not formalized.

To see this, name each live value of a frame, rather than each register, by
a label in $\{1,\ldots,W\}$; there are at most $W$ live values. Encode a
remainder step by the labels of its two operands and the label given to the
result, which takes over the label of a value never used again in this frame
(or a free label). Routes need no encoding, since labels follow values.
Dropping a value needs no encoding either, since only operands are decoded.
A subcomputation inherits the labels and is encoded by an opening marker, its
own steps, and a closing marker carrying the label it returns and the label
the result takes in the caller (or a mark that the result is discarded).
Subcomputations without remainder calls only copy a value, and are absorbed
as routes. A subcomputation consisting only of routes and one inner
subcomputation is replaced by that inner one: its returned value is either
the inner result, whose label it takes, or a copied caller value, in which
case the inner result is discarded. After these reductions, in the tree of
subcomputations every leaf performs a remainder step of its own, and every
node without one has at least two children. So there are at most $c$ of the
former and fewer than $c$ of the latter, at most $2c$ subcomputations, and at
most $5c$ steps and markers in all. With an end marker, each symbol has one
of four types and at most $W^3$ choices of labels, and words of length at
most $5c+1$ number at most $(8W^3)^{5c+1}$. The labels determine every
operand's provenance. Hence $\log C_P=O(\log W)$ for such programs.
\end{remark}

\section{Reductions}

\begin{theorem}[Reduction]\label{thm:reduction}
Let $P$ compute $\gcd(a,b)$ over $(\mathbb N,\rem,\eq_0)$, or decide coprimeness
over $(\mathbb N,\rem,\eq_0,\eq_1)$, for $a>b\geq2$, and let $c_P(a,b)$ be its
remainder complexity. Then for all $N$ and $c$,
\[
\#\{(a,b):2\leq b<a\leq N,\ \gcd(a,b)=1,\ c_P(a,b)\leq c\}
\;\leq\;C_P^{\,c+1}\cdot\max_{|S|\leq c}P_S(N).
\]
\end{theorem}

\begin{proof}
Take a convergent computation with $c_P(a,b)$ remainder calls, and let $S$ be
its shape. $S$ is valid on $(a,b)$, and it reaches $1$. For the gcd, the
output $1$ is a value of the computation and not an input, since
$a>b\geq2$. For coprimeness, suppose no value of the computation equals $1$.
Doubling every value preserves every remainder equation
($\rem(2x,2y)=2\rem(x,y)$), every $\eq_0$ answer, and every $\eq_1$ answer
(all false), and conditionals take the same branches. So the doubled tree is
a convergent computation at $(2a,2b)$ answering ``coprime''. That is wrong,
since $\gcd(2a,2b)\geq2$, and $P$ is correct on every convergent computation.
Hence some value is $1$. By Lemma~\ref{lem:sound} it is some $v_i$, and since
$a,b\geq2$ it is a produced value $v_{t+1}$.

Sum over the at most $C_P^{\,c+1}$ shapes of Lemma~\ref{lem:count}.
\end{proof}

Since about $3N^2/\pi^2$ pairs in the range are coprime, we get the following.

\begin{corollary}[A single sublogarithmic program needs very lucky shapes]\label{cor:lucky}
If one program $P$ has $c_P(a,b)\leq f(a)$ for all coprime $a>b\geq2$, with
$f$ nondecreasing and $f(N)=o(\log N)$, then for every $\varepsilon>0$ and all large $N$ some shape
of length at most $f(N)$ reaches $1$ on at least $N^{2-\varepsilon}$ pairs
$(a,b)\leq N$.
\end{corollary}

\begin{conjecture}[Per-shape bound, $\mathrm{PSB}(c_0,\delta)$]\label{conj:psb}
There are $c_0,\delta>0$ such that $P_S(N)\leq N^{2-\delta}$ for every large
$N$ and every shape $S$ of length at most $c_0\log N$.
\end{conjecture}

Some restriction on the length is necessary. For fixed $d$, Euclid's shape of
length $d$ first reaches $1$ on about $N(\ln N)^{d}/d!$ pairs. At lengths
proportional to $\ln N$ it does far better. By Dixon's theorem
\cite{dixon1970}, for all but $o(N^2)$ pairs with $N/2<a\leq N$ the number of
Euclidean steps is $(12\ln2/\pi^2)\ln N+O((\ln N)^{1/2+\varepsilon})$. By
pigeonhole, there is a single $d=(12\ln2/\pi^2+o(1))\ln N$ such that on
$N^{2-o(1)}$ coprime pairs the Euclidean remainders first reach $1$ after
exactly $d$ steps, and Euclid's shape of length $d$ reaches $1$ on all of them. So $\mathrm{PSB}(c_0,\delta)$ fails for
every $c_0>12\ln2/\pi^2$.

\begin{corollary}[PSB implies the Main Conjecture for programs]\label{cor:main}
Assume $\mathrm{PSB}(c_0,\delta)$. For every program $P$ as in
Theorem~\ref{thm:reduction}, put $r_P=\min\{c_0,\ \delta/(2\log C_P)\}$. Then
$c_P(a,b)\geq r_P\log a$ for all but $o(N^2)$ coprime pairs
$2\leq b<a\leq N$. In particular, $P$ makes at least $r_P\log a$ remainder
calls on infinitely many pairs, deterministic or not.
\end{corollary}

\begin{proof}
Take $c=\lfloor r_P\log N\rfloor\leq c_0\log N$. Then
$C_P^{\,c+1}\leq C_P\,N^{r_P\log C_P}\leq C_P\,N^{\delta/2}$, so
Theorem~\ref{thm:reduction} and PSB give at most
$C_P\cdot N^{\delta/2}\cdot N^{2-\delta}=O(N^{2-\delta/2})$ exceptions below
$N$. Finally $\log a\leq\log N$.
\end{proof}

\begin{corollary}[PSB implies tightness]\label{cor:tight}
Assume $\mathrm{PSB}(c_0,\delta)$ and let $0<\kappa<\delta/2$. Then the
companion's $D(a,b)$ satisfies $D(a,b)>\kappa\log N/\log\log N$ for all but
$o(N^2)$ coprime pairs $2\leq b<a\leq N$. Since all but $o(N^2)$ of these pairs
have $a>N/\log N$, so that $\log a\sim\log N$, together with the companion's
upper bound this gives $D=\Theta(\log a/\log\log a)$ on almost all coprime
pairs.
\end{corollary}

\begin{proof}
Put $d=\lfloor\kappa\log N/\log\log N\rfloor$; for large $N$,
$d\leq c_0\log N$. Shapes of length at most $d$ number at most
$(d+1)((d+1)!)^2\leq(d+2)^{2d+1}=N^{2\kappa+o(1)}$. A coprime pair with
$D(a,b)\leq d$ is one on which some such shape reaches $1$, so by PSB there
are at most $N^{2-\delta+2\kappa+o(1)}=o(N^2)$ of them.
\end{proof}

\section{Consistency with what is known}

\begin{itemize}
\item \textbf{Upper bounds.} The companion construction uses about $c^{2c}$
shapes, and every fixed-width program of \cite{fixedwidth} has a coefficient
$3/p$ with width $W=4^p+5$, that is $\Theta(1/\log W)$. By
Remark~\ref{rem:width},
Corollary~\ref{cor:main} would give a matching $\Omega(1/\log W)$ for register
programs of width $W$. So, granted PSB, the family of \cite{fixedwidth} is
optimal up to a constant factor, and the quantifier exchange
$\forall\varepsilon\,\exists P\to\exists P\,\forall\varepsilon$ fails.
\item \textbf{Pratt's algorithm.} It uses only $O(\log\log a)$ remainder
calls on successive Fibonacci numbers \cite[Corollary~6C.7]{book}. That set
is sparse, so this is consistent with the almost-all bound of
Corollary~\ref{cor:main}.
\item \textbf{The homomorphism method.} Corollary~\ref{cor:main} is not a
bound for all uniform processes, since uniform processes are not limited to
$C^c$ shapes. So it does not conflict with the companion's limit on the
homomorphism method.
\item \textbf{The stack.} Recursion can keep older values alive, which
suggested that one program might emulate the growing store of the companion
construction. Lemma~\ref{lem:count} shows this cannot multiply the number of
shapes beyond $C_P^{\,c+1}$: a call returns one value, and everything else in
its frame is lost.
\end{itemize}

\section{Status}
Lemma~\ref{lem:count}, Theorem~\ref{thm:reduction} and
Corollary~\ref{cor:lucky} are unconditional. Corollaries~\ref{cor:main}
and~\ref{cor:tight} depend on Conjecture~\ref{conj:psb}, which is open.
The lemma is proved for the explicit operational model of
Section~\ref{sec:model}: two sorts, no $\iota$-constants, call-by-value calls,
lazy conditionals, nondeterministic choice of equations, convergent
computations as finite trees, and cost equal to the number of $\rem$ nodes.
That general form is also formalized in Lean
(\texttt{ShortGCD/GeneralPrograms.lean}), and all modules named below are part
of the repository's normal \texttt{lake build} and axiom audit.
\begin{itemize}
\item \texttt{compile}: every evaluation of a program with arities at most
$K\geq2$ compiles into a register run of width $2K+2$ with the same number of
remainder calls. Each subterm is evaluated in a subcomputation on a copy of
the registers, so the register-program argument below applies.
\item \texttt{eval\_shape} and \texttt{card\_shapeSet\_le\_pow}: every
evaluation on inputs $a,b\geq2$ returning $1$ with at most $c$ calls has a
shape of length at most $c$ reaching $1$, drawn from a fixed set of at most
$(7A_{2K+2}^5)^{c+1}$ shapes, with $A_W$ as below. This is the Shape-Count
Lemma with $C_P=7A_{2K+2}^5$, which depends only on $K$.
\item \texttt{eval\_shape\_val}: the provenance of the output.
\item \texttt{card\_eval\_le} and \texttt{eventually\_card\_eval\_le}: the
counting consequences.
\item \texttt{double\_or\_one}, \texttt{decide\_run},
\texttt{card\_decide\_le} and \texttt{eventually\_card\_decide\_le}: the
doubling argument and the finite and asymptotic counting for programs that
decide coprimeness. Correctness is assumed only on inputs $a>b\geq2$, as in
Theorem~\ref{thm:reduction}; there the book's relation agrees with
$\gcd(a,b)=1$.
\end{itemize}
The formal model is the one of Section~\ref{sec:model}, except that values are
tagged, and both the set of function symbols and the set of bodies may be
arbitrary rather than finite. So ill-sorted terms simply have no evaluation,
and the constant depends only on the maximal arity. The formal constant is of
order $(2K+2)^{10K+10}$ rather than the finite-alphabet constant of
Lemma~\ref{lem:count}. That is harmless for the Main Conjecture, but too weak
for the optimality remark about register programs.

The register-program case is formalized in Lean, in
\texttt{ShortGCD/ShapeCount.lean}, for the width-$W$ run semantics
\texttt{RegisterProgram.Run} of \cite{fixedwidth}.
\begin{itemize}
\item \texttt{card\_run\_le} is the finite counting statement. Suppose every
shape of length at most $c$ reaches $1$ on at most $B$ pairs
$2\le b<a\le N$. Then at most $(5c+2)\,A_W^{\,5c+1}B$ such pairs admit a run
returning $1$ with at most $c$ remainder calls. Here $A_W=W^W+W^3+W^2+W+1$ is
the size of a fixed instruction alphabet.
\item \texttt{card\_accepts\_le} gives the same bound for the full verified
program.
\item \texttt{eventually\_card\_run\_le} is the asymptotic form.
$\mathrm{PSB}(c_0,\delta)$, stated with natural logarithms, gives for each
width an $r>0$ with at most $N^{2-\delta/2}$ such pairs for all large $N$ when
the budget is $r\ln N$ calls.
\end{itemize}
The proof compiles a run into a linear code of at most $5c+1$ instructions.
Routes are merged and zero-cost subcomputations are absorbed. The shape is then
read off by symbolic execution of the code. The axiom audit lists only
\texttt{propext}, \texttt{Classical.choice} and \texttt{Quot.sound}.

The formalized constant gives $\log C=O(W\log W)$, since the alphabet includes
all $W^W$ routes. The sharper $O(\log W)$ of Remark~\ref{rem:width} is not
formalized.

Corollary~\ref{cor:tight} is formalized in
\texttt{ShortGCD/TightnessReduction.lean}. \texttt{card\_within\_le\_shapes}
bounds the number of pairs with a derivation of $1$ of length at most $d$ by
$(d+1)(d+2)^{2d}\max_{|S|\leq d}P_S(N)$, and
\texttt{eventually\_card\_within\_le} states the conclusion with natural
logarithms: if every shape of length at most $\kappa\ln N/\ln\ln N$ reaches
$1$ on at most $N^{2-2\kappa-\delta}$ pairs, then eventually at most
$N^{2-\delta/2}$ pairs have a derivation of $1$ of length at most
$\kappa\ln N/\ln\ln N$.
\texttt{AlmostAll.within\_iff\_D\_le} identifies derivations of $1$ with
$D$ on coprime pairs.

\begin{thebibliography}{4}
\bibitem{book}
Yiannis N. Moschovakis.
\emph{Abstract Recursion and Intrinsic Complexity}.
Lecture Notes in Logic 48, Cambridge University Press / ASL, 2019.
\href{https://doi.org/10.1017/9781108234238}{doi:10.1017/9781108234238}.
\bibitem{companion}
Jason Davies.
\emph{The Euclidean algorithm is not optimal among uniform processes}.
Companion repository manuscript, October 2026.
\bibitem{dixon1970}
John D. Dixon.
The number of steps in the Euclidean algorithm.
\emph{Journal of Number Theory} 2(4):414--422, 1970.
\bibitem{fixedwidth}
Jason Davies.
\emph{Fixed-arity nondeterministic remainder programs: an intermediate result}.
Research note in the companion repository, \texttt{research/nondeterministic-prefix.tex}.
\end{thebibliography}
\end{document}
